\chapter{History of the Abstract Universe Project}
\label{chap:history}

\section*{Introduction}

This chapter follows a journal-style structure, tracing the project in chronological order year by year.
It seeks to illuminate how the framework originated and evolved, from its initial spark to its current state.

Working on this project has brought both immense satisfaction and deep frustration, operating in predictable cycles,
where the latter inevitably follows the former.
Perhaps the central discovery of the entire journey is simply an appreciation for the sheer, staggering complexity of everything.


\section*{History}

\begin{description}

\item[\textbf{1988}: Pain]
  The "tooth ache" cache incident ignites long discussions with colleagues on whether pain could ever be implemented as software.

\item[\textbf{1989}: Abstract Nature of Everything]
  First simulated human thought experiment: discovering that simulated humans would also discover they can duplicate their structure in a Turing machine. We did!

  Conclusion: the simulation software is recursive. There is no single additional variable between two recursion levels that would make the simulated virtual human any less real (or us any more real). Since we know the simulated human is virtual, so must we be.

  Vesa's discovery: by observing such a simulation, we would notice that simulated humans start extracting information from their spatial surroundings; they are building a Turing machine to simulate themselves!

  The universe must be abstract by nature.

\item[\textbf{1989}: The Source of Pain Mystery]
  If millions of lines of code create a suffering virtual human, what is the exact source of that suffering? Is it when the CPU moves data from memory to an internal register, or when it flips a bit?

  Realization: any computer can be implemented as a mechanical device using wooden spheres passing through wooden gates. If a vast network of wooden components creates a suffering virtual human, what does a single moving wooden component create?

  Is this \textbf{real} $\rightarrow$ \textbf{virtual} action beyond known physics?

\item[\textbf{1991}: Major Breakthrough: Human as an Axiomatic System]
  The discovery that DNA can be digitized and simulated in a digital system. Two observational assumptions serve as axioms: DNA contains the instructions on how to build pain, and no magic plays a role in its operation, it consists of ordinary particles operating according to the laws of physics.

  The laws of physics are based on mathematics, and mathematics is an axiomatic system. A human must therefore be entirely an axiomatic system. A Turing machine is an implementation of an axiomatic system; therefore, a human can be simulated.

\item[\textbf{1992}: Pain Hypothesis]
  Discovery that pain has measurable effects. A human in pain is physically distinct from a human without pain. The influence of pain can be measured just like the effects of gravity.

  Conclusion: simulated humans must feel pain, or there would be a mathematical contradiction, two identical axiomatic systems would behave differently. $2+2$ with bananas would not equal $2+2$ with apples.


\item[\textbf{1994} -- Theory of Morality and God]

  Viewing the human as an axiomatic system raises an unsettling question regarding free will and God.

  Thesis of Morality: what we define as ``right'' aligns precisely with the typical evolutionary behaviors of group-living social animals.  

  Thesis of God: we believe in God because we are social animals whose survival depends on group organization. Humans utilize the concept of leadership to maintain order, and God represents the ultimate, idealized leader. Evolution has burned
  us tendency to follow good leaders, and God represents the ultimate, idealized leader.
  
  
\item[\textbf{1993}: The Method]
  We do not know the true nature of time or a black hole singularity. However, we can simulate them. Whatever they are, their functional equivalents must exist inside the computer running the simulation. Therefore, we can study the computer, whose internal structure we perfectly understand.

\item[\textbf{1994}: The Two Time Thought Experiments]
  Discovery: time can be gradually optimized out by replacing simulation code with static data. A simulated human must still experience pain, otherwise there would be a code/data threshold where pain vanishes. Consciousness must emerge from static data.

  Derived thought experiment: in a multithreaded, preemptive computer system, consciousness inevitably emerges from pure white noise.

  The fundamental nature of the universe is static, timeless, random \textbf{noise}.

  \textit{How could gravity possibly work in static random noise?}

\item[\textbf{1995}: Collapsing Spacetime and the Equivalence of Singularity and Entropy]
  First gravity thought experiment monitoring a black hole simulation of collapsing spacetime and the entropy of its memory image. Discovery: entropy decreases as spacetime collapses. Although the simulation crashes before the singularity forms, we can extrapolate the entropy trend: a singularity is a zero-entropy state.

  Discovery: any attempt to interpret zero-entropy information as geometric space always maps to a single object, a geometric point.

\item[\textbf{1996}: Time-Reversed Black Hole Thought Experiment]
  Discovery that all geometric interpretations of increasing entropy yield an expanding universe.

  Hypothesis: an expanding universe and increasing entropy are two sides of the same coin.

  Inflationary expansion of the early universe: a random system rushing from zero entropy toward thermodynamic equilibrium.

\item[\textbf{1996}: Major Frustration: Emergence of Gravity]
  First actual simulations written in C/C++ demonstrate exploding, expanding dust clouds.

  The probability of emergent structure is practically zero.

  \textit{How could gravity possibly emerge within this framework?}

\item[\textbf{2001}: Increasing Entropy and Emergent Geometric Microstructures]
  Simulations mapping increasing entropy to geometric space while filtering microstructures reveal a log-normal distribution curve. The probability of emergent structure appears to follow a log-normal shape.

\item[\textbf{2005}: Box of Coins Experiment]
  Major discovery via the "coins in a box" thought experiment: the more ways coins can be arranged to describe a structure, the higher the probability of that structure emerging.

  How do we maximize the probability of emergence? By using some of the bits to describe a compression algorithm!

  A random emergent observer will most likely find themselves residing within a compressed structure.

\item[\textbf{2006}: Gravity, Finally]
  Predictable, smoothly evolving data compresses far better than chaotically jumping data: enter mass and inertia!

  We most likely find ourselves in a smoothly evolving, predictable universe because predictable information compresses best. That is precisely what the laws of physics represent: predictability.

\item[\textbf{2009}: Complex-Valued Wavefunction as Compression]
  Discovery that waves permeate the microcosm everywhere, we are observing compressed structure. The quantum mechanical complex-valued wavefunction is the universe's compression signature; the Born rule acts as its dithering system.

  Thought experiment: render all possible videos of yourself in all possible formats, both compressed and uncompressed. Pick one at random. What do you observe? Most likely a dithered, compressed video, because bits can be arranged to render you in vastly more ways through compression and dithering.

\item[\textbf{2010}: Black Hole Simulation]
  First serious simulation proof-of-concept demonstrating the vanishing entropy and geometric singularity hypothesis, accompanied by a semi-formal paper draft.

  Conclusions: a spacetime singularity is a zero-information state. No information, no probability of emergent particles, no spacetime curvature.

\item[\textbf{2011}: First Simulation of Emergent Gravity]
  First attempts to develop minimal toy simulations showing how physical laws can emerge simply by randomly selecting execution frames representing an evolving minimal structure (e.g., a Gaussian blob).

  Conclusion: the most compressible state histories of a Gaussian blob are the most numerous, demonstrating emergent inertia.

\item[\textbf{2012}: Emergence of Spacetime]
  First simulations demonstrating the time-reverse of a black hole: the emergence of spacetime, accompanied by a paper draft.

\item[\textbf{2014}: The Book: \textit{Good Heavens Gods? Hell No Dogs!}]
  First draft of the manuscript takes shape after an intense two-year effort.

  What's a good title for a book like this? \textit{Good Heavens Gods? Hell No Dogs!}

\item[\textbf{2018}: Why Geometry? Why Not High-Dimensional Spectral?]
  Realization: we perceive life in a geometric $3+1$ spacetime. Why geometric? Why not a high-dimensional complex-valued spectral representation in Hilbert space?

  We require well-defined inside/outside boundaries to prevent our information from leaking and mixing with the environment. Any random, extraneous information injected outside degrades compression, leading to an exponentially vanishing probability of emergence.

  Idea: General Relativity (GR) might also be about compression, a minimal description required to define informational boundaries.

  Perhaps GR and Quantum Mechanics (QM) can be unified through compression!

\item[\textbf{2019}: The Consciousness Breakthrough]
  Discovery: different optimizations of a DNA simulation result in drastically different execution traces. A simulated human is not a single specific bit sequence, but an entire equivalence class of them. This is the fundamental definition of entropy.

  If simulated consciousness emerges from a sequence of bits, we face a dilemma: either the initial axioms about DNA are wrong, or everything we observe, including pain, must emerge from bit sequences.

  Breakthrough: in a static, timeless universe, there is no external "interpreter," and raw bits do not dictate how they should be decoded (e.g., two's complement). Only a vanishingly small fraction of possible interpretations map to a Turing machine.

  Pain has no room to operate within the strict confines of a Turing machine subset. However, pain has immense room to operate \textit{outside} that subset.

  Pain may be a heuristic navigation system operating on a vast "interpretation space," where strict algorithmic methods would be far too inefficient.

\item[\textbf{2022--2023}: From DocBook to LaTeX]
  The rise of LLMs and modern AI felt like a godsend.

  Proofreading and refining texts with AI assistance to reach publishable quality.

  Converting legacy C/C++ simulations into Python, which AI handles effortlessly.

\item[\textbf{2024 -- Present}]
  Migrated all materials to GitHub.

  Adding citations and building new simulations with AI assistance.

  It hasn’t been all joy: AI proofreading gradually shifted my text into something unrecognizable, stripping away my natural tone.

  Refining and finalizing all materials will likely take several more years.

  Furthermore, thanks to AI search capabilities, it became clear that many of my foundational discoveries had already been explored by others, often long ago.

  Maybe it is time to reconsider the path forward.

\end{description}


